\documentclass[aspectratio=169, 14pt]{beamer}
\usepackage{stampfly_slides}

\lesson{10}{ws:: API リファレンス}{ws:: API Reference}

\begin{document}

% ------------------------------------------------------------------
\begin{frame}
  \titlepage
\end{frame}

% ------------------------------------------------------------------
\begin{frame}{今日のゴール / Today's Goal}
  \begin{block}{目標}
    ws:: API の全関数を把握し、独自ファームウェアで自由にセンサ・モータを操作する
  \end{block}

  \vspace{0.5em}

  \begin{itemize}
    \item ws:: API 全 33 関数をカテゴリ別に理解
    \item 各 API の使い方をコード例で確認
    \item StampFlyState による上級アクセスを知る
    \item \api{sf app new} で独自プロジェクトを作成
  \end{itemize}
\end{frame}

% ==================================================================
%   ws:: API Category Pages
% ==================================================================

% ------------------------------------------------------------------
\begin{frame}{ws:: API 概要 / API Overview}
  \centering
  {\small
  \begin{tabular}{llr}
    \hline
    \textbf{カテゴリ} & \textbf{概要} & \textbf{関数数} \\
    \hline
    Motor Control   & モータ制御・Arm/Disarm    & 7 \\
    RC Input        & コントローラスティック入力 & 4 \\
    RC Buttons      & ボタン・飛行モード         & 5 \\
    IMU Sensor      & ジャイロ・加速度センサ     & 6 \\
    Estimation      & ESKF 姿勢・高度推定        & 4 \\
    LED             & LED 色制御                 & 3 \\
    Utility         & 時刻・電圧・通信・出力     & 4 \\
    \hline
    \multicolumn{2}{l}{\textbf{合計}} & \textbf{33} \\
    \hline
  \end{tabular}
  }

  \vspace{0.5em}
  {\footnotesize
    \texttt{\#include "workshop\_api.hpp"} で全関数を使用可能。
    全関数は \texttt{ws::} 名前空間に属する。\\
    より深いアクセスには \textbf{StampFlyState} を使用（後半で解説）。
  }
\end{frame}

% ------------------------------------------------------------------
\begin{frame}[fragile]{Motor Control API（1/2）: モータ制御}
  \resizebox{\linewidth}{!}{%
  \begin{tabular}{lll}
    \hline
    \textbf{関数} & \textbf{引数} & \textbf{説明} \\
    \hline
    \texttt{motor\_set\_duty(id, duty)} & id: 1--4, duty: 0.0--1.0 & 個別モータのデューティ設定 \\
    \texttt{motor\_set\_all(duty)}      & duty: 0.0--1.0           & 全モータを同一デューティに \\
    \texttt{motor\_stop\_all()}         & ---                       & 全モータ即時停止 \\
    \texttt{motor\_mixer(t, r, p, y)}   & 各 float                 & スラスト＋姿勢ミキシング出力 \\
    \hline
  \end{tabular}
  }

  \vspace{0.3em}

  \begin{lstlisting}[style=sfcpp, basicstyle=\ttfamily\scriptsize]
// Individual motor control
ws::motor_set_duty(1, 0.3f);  // Motor 1 (FR) = 30%
ws::motor_stop_all();          // Emergency stop

// Mixer: thrust + attitude corrections
ws::motor_mixer(0.5f, 0.0f, 0.0f, 0.0f);  // Hover
  \end{lstlisting}
\end{frame}

% ------------------------------------------------------------------
\begin{frame}[fragile]{Motor Control API（2/2）: Arm / Disarm}
  \resizebox{\linewidth}{!}{%
  \begin{tabular}{lll}
    \hline
    \textbf{関数} & \textbf{戻り値} & \textbf{説明} \\
    \hline
    \texttt{arm()}       & void & モータ出力を有効化（Arm） \\
    \texttt{disarm()}    & void & モータ出力を無効化（Disarm） \\
    \texttt{is\_armed()} & bool & Arm 状態を返す（\texttt{true} = Armed） \\
    \hline
  \end{tabular}
  }

  \vspace{0.3em}

  \begin{lstlisting}[style=sfcpp, basicstyle=\ttfamily\scriptsize]
if (ws::rc_throttle() < 0.05f) ws::arm();

if (ws::is_armed()) {
    ws::motor_mixer(ws::rc_throttle(),
        ws::rc_roll(), ws::rc_pitch(), ws::rc_yaw());
} else {
    ws::motor_stop_all();
}
  \end{lstlisting}

  \vspace{0.2em}
  {\footnotesize \warn{注意:} \api{arm()} 前はモータコマンドを送っても回転しない}
\end{frame}

% ------------------------------------------------------------------
\begin{frame}[fragile]{RC Input API: スティック入力}
  \resizebox{\linewidth}{!}{%
  \begin{tabular}{lll}
    \hline
    \textbf{関数} & \textbf{範囲} & \textbf{説明} \\
    \hline
    \texttt{rc\_throttle()} & $[0.0, 1.0]$   & スロットル（上がプラス） \\
    \texttt{rc\_roll()}     & $[-1.0, 1.0]$  & ロール（右がプラス） \\
    \texttt{rc\_pitch()}    & $[-1.0, 1.0]$  & ピッチ（前がプラス） \\
    \texttt{rc\_yaw()}      & $[-1.0, 1.0]$  & ヨー（右回転がプラス） \\
    \hline
  \end{tabular}
  }

  \vspace{0.3em}

  \begin{lstlisting}[style=sfcpp, basicstyle=\ttfamily\scriptsize]
void loop_400Hz(float dt) {
    float thr = ws::rc_throttle();  // 0.0 - 1.0
    float rol = ws::rc_roll();      // -1.0 - 1.0
    float pit = ws::rc_pitch();
    float yaw = ws::rc_yaw();
    ws::motor_mixer(thr, rol, pit, yaw);  // ACRO mode
}
  \end{lstlisting}
\end{frame}

% ------------------------------------------------------------------
\begin{frame}[fragile]{RC Buttons / Mode API}
  \resizebox{\linewidth}{!}{%
  \begin{tabular}{lll}
    \hline
    \textbf{関数} & \textbf{戻り値} & \textbf{説明} \\
    \hline
    \texttt{rc\_throttle\_yaw\_button()}  & bool & 左スティック押し込みボタン \\
    \texttt{rc\_roll\_pitch\_button()}    & bool & 右スティック押し込みボタン \\
    \texttt{rc\_stabilize\_acro\_mode()}  & bool & ACRO モード判定（\texttt{true} = ACRO） \\
    \texttt{rc\_alt\_mode()}             & bool & 高度保持モード判定 \\
    \texttt{rc\_pos\_mode()}             & bool & 位置保持モード判定 \\
    \hline
  \end{tabular}
  }

  \vspace{0.3em}

  \begin{lstlisting}[style=sfcpp, basicstyle=\ttfamily\scriptsize]
void loop_400Hz(float dt) {
    if (ws::rc_stabilize_acro_mode()) {
        acro_control(dt);       // ACRO: rate control
    } else {
        stabilize_control(dt);  // Stabilize: angle control
    }
}
  \end{lstlisting}
\end{frame}

% ------------------------------------------------------------------
\begin{frame}[fragile]{IMU Sensor API}
  \resizebox{\linewidth}{!}{%
  \begin{tabular}{lll}
    \hline
    \textbf{関数} & \textbf{単位} & \textbf{説明} \\
    \hline
    \texttt{gyro\_x()}, \texttt{gyro\_y()}, \texttt{gyro\_z()}    & rad/s   & ロール/ピッチ/ヨー軸角速度（BMI270） \\
    \texttt{accel\_x()}, \texttt{accel\_y()}, \texttt{accel\_z()} & m/s$^2$ & 前方/右方/下方加速度（NED 座標系） \\
    \hline
  \end{tabular}
  }

  \vspace{0.3em}

  \begin{lstlisting}[style=sfcpp, basicstyle=\ttfamily\scriptsize]
void loop_400Hz(float dt) {
    float gx = ws::gyro_x();    // [rad/s] roll rate
    float gy = ws::gyro_y();    // [rad/s] pitch rate
    float gz = ws::gyro_z();    // [rad/s] yaw rate
    float az = ws::accel_z();   // [m/s^2] ~9.81 at rest

    ws::print(">gyro_x:%.3f", gx);
    ws::print(">accel_z:%.2f", az);
}
  \end{lstlisting}

  \vspace{0.2em}
  {\footnotesize 静止時: \texttt{accel\_z()} $\approx$ 9.81（NED 座標系、下向き正）}
\end{frame}

% ------------------------------------------------------------------
\begin{frame}[fragile]{Estimation API: ESKF 推定値}
  \resizebox{\linewidth}{!}{%
  \begin{tabular}{lll}
    \hline
    \textbf{関数} & \textbf{単位} & \textbf{説明} \\
    \hline
    \texttt{estimated\_roll()}     & rad & ESKF ロール推定角 \\
    \texttt{estimated\_pitch()}    & rad & ESKF ピッチ推定角 \\
    \texttt{estimated\_yaw()}      & rad & ESKF ヨー推定角 \\
    \texttt{estimated\_altitude()} & m   & ESKF 推定高度（正 = 上） \\
    \hline
  \end{tabular}
  }

  \vspace{0.3em}

  \begin{lstlisting}[style=sfcpp, basicstyle=\ttfamily\scriptsize]
void loop_400Hz(float dt) {
    float roll  = ws::estimated_roll();     // [rad]
    float pitch = ws::estimated_pitch();
    float alt   = ws::estimated_altitude(); // [m]
    // P control: level the drone
    ws::motor_mixer(ws::rc_throttle(),
        (0.0f-roll)*0.5f, (0.0f-pitch)*0.5f, ws::rc_yaw());
}
  \end{lstlisting}

  \vspace{0.2em}
  {\footnotesize ESKF はジャイロ+加速度+気圧+ToF+磁気+光学フローを統合（L09 参照）}
\end{frame}

% ------------------------------------------------------------------
\begin{frame}[fragile]{LED Control API}
  \resizebox{\linewidth}{!}{%
  \begin{tabular}{lll}
    \hline
    \textbf{関数} & \textbf{引数/戻り値} & \textbf{説明} \\
    \hline
    \texttt{led\_color(r, g, b)}      & r,g,b: 0--255 (uint8\_t) & LED 色を RGB で設定 \\
    \texttt{disable\_led\_task()}     & void                     & システム LED タスクを停止 \\
    \texttt{is\_led\_task\_disabled()} & bool                    & LED タスクが停止中か確認 \\
    \hline
  \end{tabular}
  }

  \vspace{0.3em}

  \begin{lstlisting}[style=sfcpp, basicstyle=\ttfamily\scriptsize]
void setup() {
    ws::disable_led_task();  // Take over LED control
}
void loop_400Hz(float dt) {
    float v = ws::battery_voltage();
    if (v > 3.8f)      ws::led_color(0, 255, 0);    // Green
    else if (v > 3.5f) ws::led_color(255, 165, 0);  // Orange
    else                ws::led_color(255, 0, 0);    // Red
}
  \end{lstlisting}

  \vspace{0.2em}
  {\footnotesize \warn{注意:} \api{disable\_led\_task()} を呼ばないとシステムが LED を上書きする}
\end{frame}

% ------------------------------------------------------------------
\begin{frame}[fragile]{Utility API}
  \resizebox{\linewidth}{!}{%
  \begin{tabular}{lll}
    \hline
    \textbf{関数} & \textbf{戻り値} & \textbf{説明} \\
    \hline
    \texttt{millis()}           & uint32\_t (ms) & 起動からの経過時間 \\
    \texttt{battery\_voltage()} & float (V)      & バッテリー電圧 \\
    \texttt{print(fmt, ...)}    & void           & printf 形式でデバッグ出力 \\
    \texttt{set\_channel(ch)}   & void           & WiFi チャネル設定（1, 6, 11） \\
    \hline
  \end{tabular}
  }

  \vspace{0.3em}

  \begin{lstlisting}[style=sfcpp, basicstyle=\ttfamily\scriptsize]
void setup() {
    ws::set_channel(6);
    ws::print("Lesson 10: API Reference");
}
void loop_400Hz(float dt) {
    static uint32_t last = 0;
    uint32_t now = ws::millis();
    if (now - last >= 20) {  // 50 Hz decimation
        last = now;
        ws::print(">battery:%.2f", ws::battery_voltage());
    }
}
  \end{lstlisting}

\end{frame}

% ==================================================================
%   Advanced: StampFlyState
% ==================================================================

% ------------------------------------------------------------------
\begin{frame}{StampFlyState: 全センサ + 推定値}
  \resizebox{\linewidth}{!}{%
  \scriptsize
  \begin{tabular}{lll}
    \hline
    \textbf{メソッド} & \textbf{ソース} & \textbf{取得データ} \\
    \hline
    \texttt{getIMUData(a, g)}       & BMI270 IMU      & 加速度 + ジャイロ生値 \\
    \texttt{getIMUCorrected(a, g)}  & バイアス補正      & バイアス除去済み IMU \\
    \texttt{getBaroData(alt, p)}    & BMP280 気圧      & 高度 [m], 気圧 [Pa] \\
    \texttt{getMagData(mag)}        & BMM150 磁気      & 磁気 x,y,z [$\mu$T] \\
    \texttt{getToFData(b, f)}       & VL53L3CX ToF    & 下方/前方距離 [m] \\
    \texttt{getFlowData(vx, vy)}    & PMW3901 フロー    & 速度 vx,vy [m/s] \\
    \texttt{getFlowRawData(...)}    & PMW3901 生値      & dx,dy [counts], squal \\
    \hline
    \texttt{getAttitude()}          & ESKF 推定        & roll, pitch, yaw [rad] \\
    \texttt{getPosition()}          & ESKF 推定        & x, y, z [m] \\
    \texttt{getVelocity()}          & ESKF 推定        & vx, vy, vz [m/s] \\
    \texttt{getGyroBias()}          & ESKF 推定        & バイアス [rad/s] \\
    \texttt{getAccelBias()}         & ESKF 推定        & バイアス [m/s$^2$] \\
    \hline
  \end{tabular}
  }

\end{frame}

% ------------------------------------------------------------------
\begin{frame}{ハードウェアセンサ仕様 / Hardware Sensors}
  \centering

  \resizebox{0.9\linewidth}{!}{%
  \begin{tabular}{llll}
    \hline
    \textbf{センサ} & \textbf{型番} & \textbf{サンプルレート} & \textbf{測定量} \\
    \hline
    IMU        & BMI270  & 400\,Hz  & 加速度 + ジャイロ \\
    気圧       & BMP280  & 50\,Hz   & 気圧 → 高度 \\
    磁気       & BMM150  & 10--30\,Hz & 磁気ベクトル \\
    ToF（下方） & VL53L3CX & 30\,Hz  & 対地距離（0--2\,m） \\
    ToF（前方） & VL53L3CX & 30\,Hz  & 前方距離（0--2\,m） \\
    光学フロー  & PMW3901 & 100\,Hz  & 対地速度 \\
    \hline
  \end{tabular}
  }

  \vspace{1em}

  {\footnotesize
    ws:: API は IMU のみ直接アクセス可（\api{gyro\_x()}, \api{accel\_x()} 等）。\\
    気圧・磁気・ToF・光学フローは \textbf{StampFlyState} 経由で取得。\\
    ESKF は全センサを統合して姿勢・位置・速度を推定（Lesson 9 参照）。
  }
\end{frame}

% ------------------------------------------------------------------
\begin{frame}{エコシステム概要 / Ecosystem Overview}
  \begin{columns}[t]
    \begin{column}{0.48\textwidth}
      \begin{exampleblock}{開発ツール}
        \begin{itemize}
          \item \api{sf build / flash / monitor}
          \item \api{sf cal gyro/accel/mag}
          \item \api{sf log wifi / capture / viz}
          \item \api{sf sim list / run}
        \end{itemize}
      \end{exampleblock}
    \end{column}
    \begin{column}{0.48\textwidth}
      \begin{exampleblock}{解析・応用}
        \begin{itemize}
          \item Python SDK（開発中）
          \item Jupyter Notebooks
          \item Genesis シミュレータ
          \item フライトログ解析
        \end{itemize}
      \end{exampleblock}
    \end{column}
  \end{columns}

  \vspace{0.5em}
  {\footnotesize ワークショップで学んだスキルは、エコシステムの全ツールで活用可能}
\end{frame}

% ------------------------------------------------------------------
\begin{frame}[fragile]{独自ファーム作成: \texttt{sf app new}}
  \begin{block}{\api{sf app new} でカスタムプロジェクトを生成}
    vehicle のコンポーネントを再利用しつつ、独自の main.cpp を記述
  \end{block}

  \vspace{0.3em}

  \begin{lstlisting}[style=sfbash]
# Create project
sf app new my_drone
# -> firmware/my_drone/ が生成される

# Build and flash
sf build my_drone
sf flash my_drone -m
  \end{lstlisting}

  \vspace{0.3em}
  {\footnotesize
    \textbf{仕組み:} \texttt{EXTRA\_COMPONENT\_DIRS} で vehicle/components と common を参照。\\
    IMU, 気圧, ToF, 光学フロー, モーター等すべてのコンポーネントが使える。
  }
\end{frame}

% ------------------------------------------------------------------
\begin{frame}[fragile]{Teleplot によるセンサ可視化}
  \begin{lstlisting}[style=sfcpp, basicstyle=\ttfamily\scriptsize]
auto& s = StampFlyState::getInstance();
Vec3 accel, gyro;
s.getIMUData(accel, gyro);
float baro_alt, pressure;
s.getBaroData(baro_alt, pressure);
float tof_bottom, tof_front;
s.getToFData(tof_bottom, tof_front);
auto att = s.getAttitude();     // ESKF roll,pitch,yaw
float alt = ws::estimated_altitude();

// Teleplot output (>name:value format)
ws::print(">baro_alt:%.2f", baro_alt);
ws::print(">tof_bottom:%.3f", tof_bottom);
ws::print(">eskf_alt:%.2f", alt);
ws::print(">eskf_roll:%.1f", att.x * 57.3f);
  \end{lstlisting}

  {\footnotesize 気圧高度 vs ToF vs ESKF 推定高度の比較で、センサ融合の効果を観察}
\end{frame}

% ------------------------------------------------------------------
\begin{frame}[fragile]{実習: 全センサ読み取り / Hands-on}
  \begin{lstlisting}[style=sfcpp, basicstyle=\ttfamily\scriptsize]
#include "workshop_api.hpp"
#include "stampfly_state.hpp"
void setup() { ws::print("Lesson 10: API Reference"); }
void loop_400Hz(float dt) {
    static uint32_t tick = 0; tick++;
    if (tick % 8 != 0) return;  // 50 Hz decimation
    // ws:: API for ESKF estimation
    ws::print(">eskf_roll:%.1f", ws::estimated_roll()*57.3f);
    ws::print(">eskf_alt:%.2f", ws::estimated_altitude());
    // StampFlyState for raw sensor data
    auto& s = StampFlyState::getInstance();
    float baro_alt, pressure;
    s.getBaroData(baro_alt, pressure);
    ws::print(">baro_alt:%.2f", baro_alt);
}
  \end{lstlisting}
\end{frame}

% ------------------------------------------------------------------
\begin{frame}{チェックポイント / Checkpoint}
  \begin{alertblock}{確認事項}
    \begin{itemize}
      \item[$\square$] ws:: API 全 33 関数のカテゴリと役割を把握した
      \item[$\square$] Motor / RC / IMU / Estimation の使い方を理解した
      \item[$\square$] StampFlyState で全センサ生値を取得できた
      \item[$\square$] Teleplot で複数センサのグラフを同時表示した
      \item[$\square$] \api{sf app new} でカスタムプロジェクトを作成できた
    \end{itemize}
  \end{alertblock}

  \vspace{0.5em}

  \begin{block}{次のレッスン}
    Lesson 11: Python SDK プログラム飛行
  \end{block}
\end{frame}

\end{document}
